\begin{statusdone}
UDP 청크 전송, 메타 필수 게이트, 적응형 레벨 선택까지 안정적으로 동작한다.
\end{statusdone}

\chapter{원리 및 설계 근거}

\section{전송해야 하는 두 종류의 데이터}
\begin{itemize}
  \item \textbf{레이더 raw ADC}: DCA1000EVM이 LVDS로 토해내는 원시 IQ 샘플.
        대역폭이 크고(수십\,Mbps) 연속 스트림이다. 손실되면 한 프레임의 신호처리가
        깨지므로 순서(sequence)와 타임스탬프가 중요하다.
  \item \textbf{웹캠 영상}: JPEG으로 압축한 프레임. 한 프레임이 UDP 단일 패킷보다
        크므로 \emph{청크 분할}이 필요하다.
\end{itemize}

\section{청크 분할과 헤더 설계}
JPEG 한 프레임을 \texttt{CHUNK\_SIZE = 1100\,B} 단위로 잘라 전송한다. 각 청크
앞에 8바이트 헤더 \texttt{>IHHI}(frame\_id, chunk\_id, total, ts\_ms)를 붙여,
수신측이 어떤 프레임의 몇 번째 조각인지/총 몇 조각인지/촬영 시각을 알 수 있게 한다.
레이더는 \texttt{>II}(seq, ts\_ms) 헤더를 붙여 순서·시각을 보존한다.

\begin{designbox}{왜 1100바이트인가}
이더넷 MTU(1500B) 안에 UDP/IP 헤더와 자체 헤더를 더해도 단편화(fragmentation)가
일어나지 않도록 보수적으로 잡은 페이로드 크기다. IP 단편화가 생기면 한 조각만
잃어도 전체 패킷이 버려져 손실률이 급증한다.
\end{designbox}

\section{적응형 레벨 시스템}
송신 화질/레이더 chirp 수를 묶은 ``레벨''을 사전 정의하고, 가용 대역폭에 맞춰
자동 선택한다.
\begin{itemize}
  \item \texttt{sender\_mode = 0}: 고정 레벨 사용.
  \item \texttt{sender\_mode = 1}: \texttt{iperf3}로 실제 대역폭을 측정 → 목표 비율
        (\texttt{bw\_target\_pct})만큼을 예산으로 삼아, 예산 이하의 가장 높은 레벨 선택.
\end{itemize}
각 레벨은 \{chirp 수, fps, 해상도, JPEG 품질, 예상 Mbps\}를 갖는다. Tailscale
직결에서 약 255\,Mbps가 측정되며, 그 안에서 레벨이 정해진다.

\chapter{구현 및 디렉토리 구조}

\section{송신측 — \texttt{src/transmission/sender.py}}
\begin{itemize}
  \item \texttt{select\_level()} : 대역폭 측정 → 레벨 결정.
  \item \textbf{레이더 경로}: TI CLI(\texttt{DCA1000EVM\_CLI\_*.exe})를
        \texttt{fpga → record → start\_record} 순으로 호출하고, UART로 AWR1642에
        \texttt{.cfg}(chirp/frame 설정)를 전송한다. \texttt{radar\_forward()}가
        DCA1000의 UDP(데이터 포트)를 받아 \texttt{(seq, ts)} 헤더를 붙여 데스크톱으로
        포워딩한다.
  \item \textbf{주기적 재시작}: DCA1000은 일정 패킷 후 멈추므로,
        ``Record is completed'' 감지 또는 \texttt{restart\_at\_seq} 도달 시
        레이더를 재설정/재시작한다(\texttt{restart\_radar}).
  \item \textbf{웹캠 경로}: \texttt{webcam\_send()} — DSHOW + MJPG로 캡처,
        JPEG 인코딩 후 청크 전송.
  \item \texttt{meta\_sender\_loop()} : 레이더 메타(.cfg 파라미터)를 TCP로 주기 전송.
\end{itemize}

\section{수신측 — \texttt{src/transmission/receiver.py}}
\begin{itemize}
  \item \texttt{radar\_receive()} : 헤더 분리 후 payload를 \emph{로컬 MATLAB}
        (\texttt{127.0.0.1:matlab\_port})으로 즉시 포워딩하고, raw ADC를 chirp 크기
        (\texttt{samples$\times$receivers$\times$chirps$\times$4})로 모아 \texttt{.bin}에 저장.
  \item \texttt{webcam\_receive()} : \texttt{frame\_id}별로 청크를 모아 \texttt{total}개가
        다 차면 재조립 → JPEG 디코드 → \texttt{\_latest\_frame} 갱신. 메인 루프는
        항상 \emph{최신 프레임}만 가져간다(\texttt{get\_latest\_frame}).
  \item \texttt{meta\_receive()} : 메타 TCP 수신 → \texttt{set\_chirp()}로
        \texttt{\_chirp\_ready} 이벤트를 세팅하고 \texttt{iqData\_RecordingParameters.mat}
        를 생성(MATLAB이 읽는 신호처리 파라미터).
\end{itemize}

\section{설정 파일}
\texttt{config/network.toml}(IP·포트), \texttt{config/sender.toml}(레벨·카메라·DCA·레이더
파라미터), \texttt{config/receiver.toml}(저장·파이프라인·YOLO·DynBG 설정)로 분리.

\chapter{문제 해결 과정}

\section{메타를 필수 게이트로}
\textbf{증상}: 레이더가 전혀 안 잡혀 모든 객체가 거리 추정(est)으로만 표시.\\
\textbf{원인}: 메타(.mat 신호처리 파라미터)가 도착하기 전 20초 타임아웃 후
``웹캠만''으로 진행하면, MATLAB 신호처리가 시작되지 못해 \texttt{targets.json}이
비어 있었다.\\
\textbf{해결}: 메타를 \emph{필수 선행 게이트}로 만들었다. \texttt{\_chirp\_ready}가
세팅될 때까지 블록(10초마다 안내)하고, 메타가 늦게 와도 도착하는 즉시
\texttt{radar\_live}를 시작한다. 디버그용 \texttt{--no-radar} 플래그로 웹캠 단독
실행도 가능.

\section{웹캠 청크 손실 대응}
재조립 버퍼에서 \texttt{frame\_id}가 5 이상 지난 미완성 프레임은 폐기하여, 손실된
청크 때문에 버퍼가 무한히 쌓이는 것을 방지한다. (수신 버퍼 \texttt{SO\_RCVBUF}
확대는 진단했으나 현재 미적용 — 향후 과제.)
